1. Vytvoř pro každou novou písničku jeden .tex soubor ve složce 'nové_písničky.'
2. Pojmenuj ho bez mezer a diakritiky, nejlepe názvem pisničky, připadně doplněným názvem autora (např. "batalion.tex", "batalion_divokej_bill.tex"), nebo akordem od kterého začíná("amazonka.tex", "amazonka_D.tex").
3. Vyzkoušej, že tvoje nová písnička funguje v 'Zpevnik_nastroj.tex' a oprav případné chyby ve své písničce.
4. Ulož písničku i do složky 'songs.'

(5. Pokud chceš mít u písničky obrázek, nahraj ho do složkky "images" a pod poslední sloku/refrén připoj příkaz: \includegraphics[scale=0.x]{images/"jmeno_tveho_obrazku.png"} x prostě párkrát zkus tak, aby obrázek byl "akorát" vidět, typicky stačí 0.1-0.8 (bacha, musí to být desetinná tečka, ne čárka!!! Napíšeš ji jako normální tečku.). Obrázek nemusí být .png, může být i .jpg a možná i další, prostě vyzkoušej :D. 

POZN.: U "normálních písniček" není třeba psát akordy do každé sloky, stačí nad první. Pokud je však změna nebo si to přeješ, můžeš je nadepsat i dále.

POZN. 2: Text písniček určitě NEPIŠ rukou!!! To trvá strašně dlouho, najdi si někde text a ten zkopíruj :D.

PSANÍ PÍSNIČKY
V samotném souboru bude použito několik příkazů: 
Znak '\' se píše pomocí AltGr + Q, znak '{' pomocí AltGr + B, znak '}' pomocí AltGr + N .

ZAČÁTEK A KONEC:
\song{název písničky}       - začátek písničky. Je zvykem psát \song{název písničky (název autora)}. Tím každý soubor začne.
\esong{}		            - konec písničky. Tím soubor skončí. Do závorek {} se zde nic nepíše.

SLOKA A REFRÉN:
\ve{} 			- nová (číslovaná) sloka. Do závorek {} se zde nic nepíše (z angl. Verse)
\cho{}			- refrén (odrážka, která rovnou napíše R:). Do závorek {} se zde nic nepíše (z angl. Chorus)
\re{odrážka}	- jiná odrážka. Třeba prázdná (\re{}), nebo s vlastním textem (\re{*:}, \re{R2:}, \re{holky:} (angl. Rest)

pozn.: po příkazu \song... musí následovat nějaký z těchto 3 odrážkových. např. "capo 3" tedy napíšu takhle:

	\song{název (autor)}

	\re{}
	capo 3

	\ve{}
	sloka...


AKORDY:
    \ch{akord}		- akord. Bude umístěn nad následujícím znakem.
    			- Nepíše se za ním mezera. např. "vlasy" s akordem začínajícím nad "l" budou psány jako "v\ch{akord}lasy"
    			- v závorce zde může být i více akordů za sebou, oddělených mezerou, typicky na konci řádku, kde pokračuje nějaká mezihra ("ko\ch{akord1 akord2 akord3}nec."
    			- pokud má akord začínat nad znakem s diakritikou (s háčkem nebo s čárkou), je potřeba tento znak uzavřít taky do {} závorek ("... \ch{akord}{Ř}eka" 
            
\guitarFlat{}       - napíše znak pro b (béčko), používá se uvnitř akordů
                    - např. \ch{B\guitarFlat{}7}  bude Bb7
\guitarSharp{}      - napíše znak pro # (křížek), používá se uvnitř akordů
                    -např. \ch{F\guitarSharp{}maj7} bude F#maj7

OBECNÉ:
\\                  - ukončí (zalomí) řádek.
                    -obecně NEdoporučujeme dávat na poslední řádek před novou slokou, nebo refrénem. Příkazy '\ve{}' a '\cho{}' automaticky udělají řádek mezeru. Viz ukázka níže :D.
(prázdný řádek)	    - ukončí odstavec (vloží vertikální mezeru před další odstavec
%                   - Udělá ze zbytku řádku poznámku, která se nevytiskne a neprojeví
\uv{text}		    - vloží text do (českých) uvozovek

A pak je spousta příkazů pro profíky.

UKÁZKA:
\song{Já s tebou žít nebudu (Nerez)}

\ve{}
\ch{Emi}Bylas' jak poslední hlt \ch{F\guitarSharp{}}vína,\\  
\ch{H7}sladká a k ránu bolavá,\ch{Emi Edim H7}{}\\
\ch{Emi}za nehty výčitky a špí\ch{F\guitarSharp}na,\\
\ch{H7}{ř}íkalas': dnešní noc je tvá.\ch{Emi Edim H7}{}

\cho{}
\ch{G}Co my dva z lásky vlastně \ch{H}máme,\\
\ch{Ami}hlubokou \ch{Edim}{š}achtou padá zd\ch{H7}viž, já říkám \uv{kiš-kiš},\\
\ch{Emi}navrch má vždycky těžký \ch{F\guitarSharp}kámen\\
\ch{H7}a my jsme v koncích čím dál blíž,\ch{Emi Edim H7}{}\\
/: \ch{Emi}A me tu ha nadži vava \ch{Ami}jaj dari \ch{Edim H7}dari daj. :/

\ve{}
Zejtra tě potkám za svým stínem,\\
neznámí známí v tramvaji,\\
v cukrárně kávu s harlekýnem,\\
hořká a sladká splývají.

...

\includegraphics[scale=0.x]{images/"nazev"}
\esong{}
